\paragraph{解析分歧半径、二维全息承载、可见 Jacobian 赤字与稀疏 prime-language 的共尾障碍}
在核版 RH 的临界边界、\(2\)-进全息地址、\(S_4\)-Hurwitz 商塔的 Jacobian/Prym 分解以及 Zeckendorf--G\"odel 素数语言的密度理论都已确定之后，还可把这些看似分散的接口继续压缩为下列七条终端结论。它们分别刻画 RH-sharp 阈值与解析分歧半径的严格分离、全息地址中首差深度的精确赋值律、前缀柱集与 \(2^L\)-进仿射球的严格对应、固定二维 \(2\)-进 ambient space 与有限秩加法账本不可能性之间的结构分裂、前缀恢复的正线性码率下界、商曲线可见因子对九维 Prym 三重块的严格维数赤字，以及稀疏 prime-register 语言在强局部约束下仍保持的正密度共尾性。

\begin{theorem}[解析分歧半径与 RH-sharp 阈值的严格分离]\label{thm:xi-time-part9w-analytic-branch-radius-rhsharp-separation}
记
\[
P_2(\theta):=\log \rho(e^\theta),
\qquad
R_\theta:=\text{\(P_2\) 在 \(\theta=0\) 处的 Taylor 收敛半径},
\]
并设 \(u_R>1\) 为 kernel RH/Ramanujan 条件的唯一正实临界参数，
\[
\theta_R:=\log u_R.
\]
则有严格不等式
\[
0<\theta_R<R_\theta.
\]
更具体地，记
\[
f(\lambda;u):=\lambda^3-2\lambda^2-(u+1)\lambda+u
\]
为碰撞核的控制性三次特征因子，则对每个 \(u>0\)，其判别式
\[
\Delta(u)=4u^3+25u^2+88u+8
\]
恒满足 \(\Delta(u)>0\)，并且三条实谱支始终保持固定根型
\[
\lambda_+(u)>1,\qquad 0<\lambda_0(u)<1,\qquad \lambda_-(u)<0.
\]
再记
\[
\Lambda(u):=\max\{|\lambda_0(u)|,|\lambda_-(u)|,1\}.
\]
则
\[
\textsf{RH}(u)\iff
\begin{cases}
\text{恒真},& 0<u\le 1,\\[1mm]
(-\lambda_-(u))^2\le \lambda_+(u),& u>1,
\end{cases}
\]
并且在 \(u=u_R\) 处满足
\[
\Lambda(u_R)=\lambda_+(u_R)^{1/2}.
\]
因此，RH-sharp 失效并非由分歧点逼近、复谱生成或根碰撞触发，而是发生在解析域内部的纯实谱阈值穿越。
\end{theorem}
\begin{proof}
由定理 \ref{thm:xi-time-part9pb-rh-transition-analytic-not-branch}，
\[
\theta_R\approx 1.27888224610494,
\qquad
R_\theta\approx 2.76230289346087,
\]
故
\[
0<\theta_R<R_\theta.
\]

另一方面，定理 \ref{thm:xi-time-part9pb-rh-transition-analytic-not-branch} 的证明已给出
\[
\Delta(u)=4u^3+25u^2+88u+8>0
\qquad (u>0),
\]
从而正实轴上不会发生代数分歧塌缩。再由命题
\ref{prop:real-input-40-collision-all-real}，三次因子 \(f(\lambda;u)\) 对每个 \(u>0\) 都具有三条互异实根，并按
\[
\lambda_+(u)>1,\qquad 0<\lambda_0(u)<1,\qquad \lambda_-(u)<0
\]
排序。

定理 \ref{thm:xi-time-part62d-collision-rh-negative-branch-control} 进一步说明：\(\lambda_0(u)\) 从不决定 RH 窗口边界；当 \(u>1\) 时，
\[
\Lambda(u)=|\lambda_-(u)|=-\lambda_-(u),
\]
于是
\[
\textsf{RH}(u)\iff \Lambda(u)\le \lambda_+(u)^{1/2}
\iff
(-\lambda_-(u))^2\le \lambda_+(u).
\]
特别地，在 \(u=u_R\) 处有
\[
\Lambda(u_R)=\lambda_+(u_R)^{1/2}.
\]
因此，RH-sharp 边界既不由解析边界触发，也不由复共轭谱支生成触发，而只由解析域内部的纯实谱不等式翻转触发。证毕。
\end{proof}

\begin{theorem}[\texorpdfstring{$2^L$}{2^L}-进全息地址的首差--赋值恒等式]\label{thm:xi-time-part9w-holographic-first-difference-valuation}
沿用定理 \ref{thm:killo-godel-semidir-affine-2adic} 的记号。设 \(E\) 为有限事件集，
\[
\mathrm{code}:E\hookrightarrow \NN,
\qquad
M:=\max_{e\in E}\mathrm{code}(e),
\qquad
B:=2^L>M,
\]
并取 primitive 地址向量
\[
v\in T:=T_2(E_{\min})\setminus 2T.
\]
对任意无限事件流
\[
\mathbf e=(e_1,e_2,\dots)\in E^{\NN},
\]
定义其全息地址
\[
X(\mathbf e):=\sum_{t\ge 1}\mathrm{code}(e_t)\,B^{t-1}v\in T.
\]
若 \(\mathbf e\neq \mathbf e'\)，并记其首个分歧位置为
\[
\tau(\mathbf e,\mathbf e')
:=
\min\{t\ge 1:\ e_t\neq e_t'\},
\]
则有精确赋值恒等式
\[
X(\mathbf e)-X(\mathbf e')
\in
B^{\tau(\mathbf e,\mathbf e')-1}T\setminus B^{\tau(\mathbf e,\mathbf e')}T.
\]
\leanverified{paper\_xi\_time\_part9w\_holographic\_first\_difference\_valuation}
\end{theorem}
\begin{proof}
记 \(\tau:=\tau(\mathbf e,\mathbf e')\)，并置
\[
\delta:=\mathrm{code}(e_\tau)-\mathrm{code}(e_\tau')\neq 0.
\]
由于 \(\mathrm{code}\) 取值于 \(\{0,\dots,M\}\) 且 \(M<B\)，有
\[
|\delta|<B.
\]
把差值从首个分歧项起展开，得到
\[
\begin{aligned}
X(\mathbf e)-X(\mathbf e')
&=
\sum_{t\ge \tau}
\bigl(\mathrm{code}(e_t)-\mathrm{code}(e_t')\bigr)B^{t-1}v\\
&=
B^{\tau-1}
\left(
\delta v+
\sum_{t>\tau}
\bigl(\mathrm{code}(e_t)-\mathrm{code}(e_t')\bigr)B^{t-\tau}v
\right).
\end{aligned}
\]
括号内第二项属于 \(BT\)，故
\[
X(\mathbf e)-X(\mathbf e')\in B^{\tau-1}T.
\]

现证它不属于 \(B^\tau T\)。若反设
\[
X(\mathbf e)-X(\mathbf e')\in B^\tau T,
\]
则必有
\[
\delta v+By\in BT
\]
对某个 \(y\in T\) 成立，从而 \(\delta v\in BT\)。设
\[
\delta=2^s\delta_0,
\qquad
0\le s<L,
\qquad
\delta_0\in \ZZ_2^\times.
\]
则
\[
2^s\delta_0 v\in 2^L T.
\]
由于 \(\delta_0\) 在 \(\ZZ_2\) 中为单位，可消去得
\[
v\in 2^{L-s}T\subset 2T,
\]
与 \(v\notin 2T\) 矛盾。故
\[
X(\mathbf e)-X(\mathbf e')\notin B^\tau T.
\]
于是所示精确赋值恒等式成立。证毕。
\end{proof}

\begin{corollary}[前缀柱集在全息像中的仿射球等式]\label{cor:xi-time-part9w-cylinder-affine-ball-identity}
沿用定理 \ref{thm:xi-time-part9w-holographic-first-difference-valuation} 的记号。固定长度 \(n\) 的前缀
\[
a=(a_1,\dots,a_n)\in E^n,
\]
并记
\[
x_a:=\sum_{t=1}^{n}\mathrm{code}(a_t)\,B^{t-1}v.
\]
定义 cylinder
\[
[a]
:=
\{\mathbf e\in E^{\NN}:\ e_1=a_1,\dots,e_n=a_n\}.
\]
则有精确仿射球公式
\[
X([a])
=
\bigl(x_a+B^nT\bigr)\cap X(E^{\NN}).
\]
因此，长度 \(n\) 的事件前缀并非仅仅近似决定地址，而是恰好切出一个 \(B^n\)-级 \(2\)-进仿射球。
\end{corollary}
\begin{proof}
若 \(\mathbf e\in[a]\)，则
\[
X(\mathbf e)-x_a
=
\sum_{t>n}\mathrm{code}(e_t)\,B^{t-1}v
\in B^nT,
\]
故
\[
X([a])\subseteq \bigl(x_a+B^nT\bigr)\cap X(E^{\NN}).
\]

反之，设 \(X(\mathbf e)\in (x_a+B^nT)\cap X(E^{\NN})\)。任取一条固定尾流
\[
\boldsymbol\eta=(\eta_1,\eta_2,\dots)\in E^{\NN},
\]
并令
\[
\mathbf f:=(a_1,\dots,a_n,\eta_1,\eta_2,\dots)\in[a].
\]
则 \(X(\mathbf f)-x_a\in B^nT\)，因此
\[
X(\mathbf e)-X(\mathbf f)\in B^nT.
\]
若 \(\mathbf e\) 的前 \(n\) 位中有某一位偏离 \(a\)，则 \(\mathbf e\) 与 \(\mathbf f\) 的首个分歧位置必不超过 \(n\)。由定理
\ref{thm:xi-time-part9w-holographic-first-difference-valuation}，
\[
X(\mathbf e)-X(\mathbf f)\in B^{t-1}T\setminus B^tT
\]
对某个 \(t\le n\) 成立，这与 \(X(\mathbf e)-X(\mathbf f)\in B^nT\) 矛盾。故 \(\mathbf e\) 的前 \(n\) 位必恰为 \(a\)，亦即 \(\mathbf e\in[a]\)。证毕。
\end{proof}

\begin{theorem}[有限二维全息实现与有限秩账本不可能性的严格分离]\label{thm:xi-time-part9w-twoadic-ambient-vs-additive-ledger-separation}
沿用定理 \ref{thm:killo-godel-semidir-affine-2adic} 的记号。任意 canonical 有限时深 prime-register 子单子
\[
H_{\mathrm{can}}\subset \mathcal R\rtimes_\Sigma \NN
\]
都可忠实实现为固定二维 \(2\)-进 Tate 模块
\[
T_2(E_{\min})\cong \ZZ_2^2
\]
上的仿射变换单子。另一方面，正整数乘法幺半群 \(\NN_{\times}\) 不能单射同态嵌入任何有限生成交换可消去加法幺半群，因而也不能嵌入任何有限生成加法账本或有限生成交换群。再者，对任意有限素数集
\[
S\subset \mathbb P,
\qquad
G_S:=\ZZ[S^{-1}],
\qquad
\Sigma_S:=\widehat{G_S},
\]
局部化相位采样仅能唯一恢复该有限素数集 \(S\)，等价地仅能恢复对偶短正合列中的 profinite 核
\[
\prod_{p\in S}\ZZ_p.
\]
因此，在本文框架中，真正的障碍并不是缺乏有限维 ambient space，而是 cofinal prime-multiplicative 结构无法被 finite-rank additive ledger 保持。
\end{theorem}
\begin{proof}
定理 \ref{thm:xi-time-part75-faithful-finite-dimensional-realization-affine-necessity} 已证明：当 \(B=2^L>M\) 时，存在忠实单子嵌入
\[
\Psi:H_{\mathrm{can}}\hookrightarrow \mathrm{Aff}\!\bigl(T_2(E_{\min})\bigr),
\qquad
T_2(E_{\min})\cong \ZZ_2^2.
\]
故任意有限时深的 canonical prime-register 机制都可以放入固定二维 \(2\)-进 ambient space 中。

同一定理还证明：\(\NN_{\times}\) 不能嵌入任何有限生成交换可消去幺半群，因此尤其不能嵌入 \((\NN^k,+)\) 或任何有限生成交换群的加法账本。故一切有限秩 additive bookkeeping 都不足以忠实保存 prime-multiplicative 结构。

最后，由定理 \ref{thm:xi-localized-phase-sampling-closed-form}，局部化相位采样函数 \(\Sigma^{\mathrm{ph}}_{G_S}\) 唯一恢复素数集 \(S\)，并进一步唯一恢复 Pontryagin 对偶短正合列
\[
0\to \prod_{p\in S}\ZZ_p\to \Sigma_S\to \TT\to 0
\]
中的 profinite 核。由于 \(S\) 有限，任何固定 prime-stage 只能携带有限多个 prime 方向，不可能承载全体素数的共尾结构。综上得到所述严格分离。证毕。
\end{proof}

\begin{corollary}[\texorpdfstring{$2^L$}{2^L}-进前缀恢复的最小码率下界]\label{cor:xi-time-part9w-prefix-recovery-minimal-rate}
沿用定理 \ref{thm:xi-time-part9w-holographic-first-difference-valuation} 的全息地址记号，并沿用定理 \ref{thm:xi-optimal-reconstruction-success-probability-constant-bias} 的微态恢复口径。若在分辨率 \(m\) 上，一种恢复协议只允许读取全息地址 \(X(\mathbf e)\) 的前 \(n\) 个 \(B\)-进位，并要求在均匀微态输入下达到成功率
\[
\mathrm{Succ}_m\ge 1-\varepsilon,
\qquad
0<\varepsilon<1,
\]
则必有
\[
n\log_2 B
\ge
m-\log_2|X_m|+\log_2(1-\varepsilon).
\]
特别地，由 \(|X_m|=F_{m+2}\sim \varphi^m\) 可得
\[
\liminf_{m\to\infty}\frac{n}{m}
\ge
\frac{\log_2(2/\varphi)}{\log_2 B}.
\]
\end{corollary}
\begin{proof}
读取前 \(n\) 个 \(B\)-进位，至多只能区分
\[
B^n
\]
种不同前缀标签；换言之，该协议所使用的附加字母表大小 \(M_m\) 至多满足
\[
M_m\le B^n.
\]
因此，其可用辅助信息量
\[
L_m:=\log_2 M_m
\]
满足
\[
L_m\le n\log_2 B.
\]

另一方面，推论 \ref{cor:xi-reconstruction-cardinality-strong-converse} 给出：任何在均匀输入下达到成功率 \(\delta\in(0,1)\) 的恢复协议都必须满足
\[
L_m\ge m-\log_2|X_m|+\log_2\delta.
\]
取 \(\delta=1-\varepsilon\)，便得到
\[
n\log_2 B
\ge
L_m
\ge
m-\log_2|X_m|+\log_2(1-\varepsilon),
\]
即第一式。

再由 \(|X_m|=F_{m+2}\) 与 Binet 渐近
\[
\log_2 F_{m+2}=m\log_2\varphi+O(1),
\]
两边同除以 \(m\log_2 B\) 并取下极限，即得
\[
\liminf_{m\to\infty}\frac{n}{m}
\ge
\frac{\log_2(2/\varphi)}{\log_2 B}.
\]
证毕。
\end{proof}

\begin{theorem}[商曲线可见因子的维数赤字与 Prym 严格多数]\label{thm:xi-time-part9w-visible-jacobian-deficit-prym-majority}
沿用定理 \ref{thm:killo-s4-genus49-jacobian-computable-isogeny} 与结论 \ref{con:xi-time-part9n-s4-all-subgroup-genera} 的记号。记
\[
J_{\mathrm{vis}}:=J(H)\times J(Z)^2\times J(C_F)^3,
\qquad
J_{\mathrm{hid}}:=P^3,
\]
其中
\[
g(H)=5,\qquad g(Z)=4,\qquad g(C_F)=3,\qquad \dim P=9.
\]
则
\[
\dim J_{\mathrm{vis}}=5+2\cdot 4+3\cdot 3=22,
\qquad
\dim J_{\mathrm{hid}}=3\cdot 9=27.
\]
特别地，
\[
27>22.
\]
因此，九维 Prym 因子的三重块不仅非平凡，而且在维数上严格超过全部商曲线可见因子的总和。于是，任何只通过商曲线
\[
X/A_4,\qquad X/D_8,\qquad X/S_3
\]
所对应的可见 Jacobian 乘积来读取证书的阿贝尔读出函子，都必然遗漏 \(J(X)\) 的严格多数部分。
\end{theorem}
\begin{proof}
由定理 \ref{thm:killo-s4-genus49-jacobian-computable-isogeny}，
\[
J(X)\sim_{\QQ}J(H)\times J(Z)^2\times J(C_F)^3\times P^3.
\]
再由结论 \ref{con:xi-time-part9n-s4-all-subgroup-genera} 与相关后续结论，
\[
g(H)=5,\qquad g(Z)=4,\qquad g(C_F)=3,\qquad \dim P=9.
\]
对 Jacobian 有 \(\dim J(C)=g(C)\)，故
\[
\dim J_{\mathrm{vis}}=5+2\cdot 4+3\cdot 3=22,
\qquad
\dim J_{\mathrm{hid}}=3\cdot 9=27.
\]
于是
\[
22+27=49=\dim J(X),
\]
并且 \(27>22\)。

此外，结论 \ref{con:xi-time-part9n-s4-cyclic-quotient-genus-spectrum} 已给出
\[
H^0(X,\Omega_X^1)\cong
5\cdot \mathrm{sgn}\oplus 4\cdot V_2\oplus 3\cdot V_3\oplus 9\cdot V_3',
\]
其中平凡表示不出现。故缺失的 \(27\) 维部分并不是某个可由不变平凡方向补回的低维尾项，而是几何上真正主导的 Prym 扇区。凡是仅经由 \(J_{\mathrm{vis}}\) 因子化的阿贝尔读出，都无法探测 \(P^3\) 这一严格多数块。证毕。
\end{proof}

\begin{theorem}[稀疏 prime-register 语言的正密度共尾性]\label{thm:xi-time-part9w-zg-positive-density-cofinal-tail}
定义
\[
\mathrm{ZG}
:=
\left\{
\prod_{k\ge 1}p_k^{z_k}:\ 
z_k\in\{0,1\},\ 
z_kz_{k+1}=0,\ 
z_k=0\ \text{对充分大 }k
\right\},
\]
即平方自由且不含相邻素因子的 Zeckendorf--G\"odel 素数语言。则：
\begin{enumerate}
  \item \(\mathrm{ZG}\) 具有正的自然密度
  \[
  \delta(\mathrm{ZG})\in(0,1);
  \]
  \item 对任意有限素数集 \(S\subset\mathbb P\)，都存在 \(n\in \mathrm{ZG}\) 使其素因子支撑满足
  \[
  \operatorname{supp}(n)\subset \mathbb P\setminus S;
  \]
  \item 因而，“禁止相邻素因子”这一强局部稀疏约束并不会把 prime-register 语言压成可忽略尾集；它仍保持正密度与共尾素数支撑。
\end{enumerate}
\end{theorem}
\begin{proof}
第一条由定理 \ref{thm:xi-zg-abel-residue-log-density} 直接得到：该定理已证明
\[
\delta_{\mathrm{ZG}}=\delta(\mathrm{ZG})\in(0,1).
\]

对第二条，任取有限素数集 \(S\subset\mathbb P\)。由于 \(\mathbb P\setminus S\) 非空，可取某个素数 \(p_j\notin S\)。令
\[
z_j=1,
\qquad
z_k=0\quad (k\neq j).
\]
则对应整数正是
\[
n=p_j\in \mathrm{ZG},
\]
因为它显然平方自由，且仅激活一个素数指标，不可能违反相邻约束；同时
\[
\operatorname{supp}(n)=\{p_j\}\subset \mathbb P\setminus S.
\]
故第二条成立。

第三条只是前两条的直接合并。再结合定理 \ref{thm:xi-localized-phase-sampling-closed-form}，任意固定有限 prime-stage 只恢复其自身有限素数支撑 \(S\)，因而不可能穷尽 \(\mathrm{ZG}\) 的共尾支撑。证毕。
\end{proof}

\paragraph{有限维 Lie 相位读出、dyadic solenoid 圆因子与零温四点骨架}
在整数共振梯的 Shannon 非可寻址性、有限二维 \(2\)-进 ambient space 的忠实实现分裂以及真实输入 \(40\) 状态核的镜像慢模已经确定之后，还可进一步抽取出下列五条后果。它们分别把有限维紧 Lie 相位上的 \(L^2\) Fourier 读出障碍、dyadic solenoid 圆因子的绝对连续排斥、显式原子 Witt 因子剥离下的慢混合主阶不变性，以及零温归一化谱中的四点骨架与双层分裂统一到同一条相位--谱层级链上。

\begin{theorem}[有限维 Lie 相位的 \texorpdfstring{$L^2$}{L2} 共振读出障碍]\label{thm:xi-time-part9w-lie-phase-l2-resonance-readout-obstruction}
设 \(G\) 为有限维紧阿贝尔 Lie 群，\(\widehat G\) 为其角色格，\(\gamma\in\widehat G\) 为无限阶角色。则不存在 \(f\in L^2(G)\) 使
\[
\widehat f(u\gamma)=C_\varphi(u)\qquad(u\in\ZZ).
\]
因而黄金 bulk 共振梯不可能是任何有限维紧 Lie 相位空间上有限能量可观测量的单频 Fourier 读出。
\end{theorem}
\begin{proof}
由 Plancherel 定理，
\[
\sum_{\chi\in\widehat G}|\widehat f(\chi)|^2=\|f\|_{L^2(G)}^2<\infty.
\]
由于 \(\gamma\) 具有无限阶，映射 \(u\mapsto u\gamma\) 在 \(\ZZ\) 上单射，从而
\[
\sum_{u\in\ZZ}|\widehat f(u\gamma)|^2\le \sum_{\chi\in\widehat G}|\widehat f(\chi)|^2<\infty.
\]
另一方面，定理 \ref{thm:fold-sigma-phi-diverges} 已给出
\[
\sum_{u\in\ZZ}C_\varphi(u)^2=+\infty.
\]
故若 \(\widehat f(u\gamma)=C_\varphi(u)\)，便与上式矛盾。证毕。
\end{proof}

\begin{corollary}[dyadic solenoid 圆因子的绝对连续障碍]\label{cor:xi-time-part9w-dyadic-solenoid-circle-factor-nonac}
记 \(\Sigma_{\{2\}}:=\widehat{\ZZ[1/2]}\)，并令
\[
0\longrightarrow \ZZ_2\longrightarrow \Sigma_{\{2\}}
\xrightarrow{\pi_{\{2\}}}\TT\longrightarrow 0
\]
为其标准圆商短正合列。则不存在 \(\Sigma_{\{2\}}\) 上的 Borel 概率测度 \(\nu\)，使得
\[
\eta:=\pi_{\{2\},*}\nu
\]
同时满足
\[
\widehat\eta(n)=C_\varphi(n)\qquad(n\in\ZZ)
\]
与 \(\eta\ll m_{\TT}\)。等价地，这样的圆因子测度不能写成任何 \(L^1(\TT)\) 密度；并且它也不可能来自任何 \(L^2(\TT)\) 密度。
\end{corollary}
\begin{proof}
若 \(\eta\ll m_{\TT}\)，则在有限测度空间 \(\TT\) 上可写成 \(d\eta=h\,dm_{\TT}\) 且 \(h\in L^1(\TT)\)。由 Riemann--Lebesgue 引理，
\[
\widehat\eta(n)\longrightarrow 0 \qquad (|n|\to\infty).
\]
但推论 \ref{cor:fold-resonance-ladder-fib-luc-limits} 与推论 \ref{cor:fold-resonance-ladder-no-decay-no-limit} 表明，\(C_\varphi(n)\) 沿 Fibonacci 与 Lucas 子列分别趋于两个不同的正常数，尤其 \(\limsup_{|n|\to\infty}C_\varphi(n)>0\)。故 \(\widehat\eta(n)=C_\varphi(n)\) 不可能成立。
若进一步有 \(h\in L^2(\TT)\)，则 Parseval 定理给出 \((\widehat\eta(n))_{n\in\ZZ}\in\ell^2(\ZZ)\)；而定理 \ref{thm:fold-sigma-phi-diverges} 表明 \((C_\varphi(n))_{n\in\ZZ}\notin\ell^2(\ZZ)\)。故 \(L^2\) 密度同样不可能。证毕。
\end{proof}

\begin{theorem}[原子 Witt 因子剥离下的慢混合主阶不变性]\label{thm:xi-time-part9w-atomic-witt-surgery-slowmixing-invariance}
沿用推论 \ref{cor:real-input-40-factor-sqrtu-eigs}、命题 \ref{prop:real-input-40-output-mirror-branch} 与定理 \ref{thm:xi-output-potential-endpoint-near-degeneracy-sqrtu-relaxation} 的记号。设
\[
\Delta(z,u)=\det(I-zM(u))=(1-uz^2)F(z,u),
\qquad
\zeta_{\mathrm{core}}(z,u):=F(z,u)^{-1}.
\]
记 \(\lambda_1(u)\) 为 Perron 根，\(\Lambda_2(u)\) 为全谱最大非主模特征值，\(\lambda_{\mathrm{mir}}(u)<0\) 为核心块的镜像慢模，并定义
\[
\tau_{\mathrm{full}}(u):=\left[-\log\!\left(\frac{\Lambda_2(u)}{\lambda_1(u)}\right)\right]^{-1},
\qquad
\tau_{\mathrm{core}}(u):=\left[-\log\!\left(\frac{|\lambda_{\mathrm{mir}}(u)|}{\lambda_1(u)}\right)\right]^{-1}.
\]
则
\[
\frac{\tau_{\mathrm{full}}(u)}{\tau_{\mathrm{core}}(u)}\longrightarrow 1
\qquad (u\to+\infty).
\]
因此剥离显式原子 Witt 因子 \((1-uz^2)^{-1}\) 以后，主导慢混合的首阶时间尺度保持不变。
\end{theorem}
\begin{proof}
由定理 \ref{thm:xi-output-potential-endpoint-near-degeneracy-sqrtu-relaxation}，
\[
1-\frac{\Lambda_2(u)}{\lambda_1(u)}=\frac{2d}{c}u^{-1/2}+O(u^{-1}).
\]
另一方面，命题 \ref{prop:real-input-40-output-mirror-branch} 给出
\[
1-\frac{|\lambda_{\mathrm{mir}}(u)|}{\lambda_1(u)}=\frac{2d}{c}u^{-1/2}+O(u^{-1}),
\qquad
\frac{\sqrt u}{\lambda_1(u)}\longrightarrow \frac1c<1.
\]
故显式原子分支 \(\pm\sqrt u\) 在归一化后始终严格位于镜像慢模之内，而
\[
-\log\!\left(\frac{\Lambda_2(u)}{\lambda_1(u)}\right)
\quad\text{与}\quad
-\log\!\left(\frac{|\lambda_{\mathrm{mir}}(u)|}{\lambda_1(u)}\right)
\]
具有相同的主阶展开 \(\frac{2d}{c}u^{-1/2}+O(u^{-1})\)。两边取倒数即得所示极限。证毕。
\end{proof}

\begin{theorem}[零温归一化谱的四点骨架]\label{thm:xi-time-part9w-zero-temp-normalized-spectrum-fourpoint-skeleton}
沿用命题 \ref{prop:real-input-40-output-mirror-branch} 的常数 \(c\)，并置
\[
b:=\frac1{c^2}\in(0,1).
\]
则 \(b\) 满足
\[
1-6b+9b^2-b^3-b^4=0.
\]
对真实输入 \(40\) 状态核，定义
\[
\mathrm{Spec}_{\mathrm{norm}}(u):=\left\{\frac{\lambda}{\lambda_1(u)}:\lambda\in\mathrm{Spec}(M(u))\right\}.
\]
若 \(\mathcal L\) 是 \(\mathrm{Spec}_{\mathrm{norm}}(u)\) 沿某子列 \(u\to+\infty\) 的 Hausdorff 极限点集，则
\[
\{1,-1,\sqrt b,-\sqrt b\}\subseteq \mathcal L.
\]
\end{theorem}
\begin{proof}
归一化 Perron 根恒等于 \(1\)，故 \(1\in\mathcal L\)。由命题 \ref{prop:real-input-40-output-mirror-branch}，
\[
\frac{\lambda_{\mathrm{mir}}(u)}{\lambda_1(u)}
=
-\left(1-\frac{2d}{c}u^{-1/2}+O(u^{-1})\right)
\longrightarrow -1,
\]
故 \(-1\in\mathcal L\)。再由推论 \ref{cor:real-input-40-factor-sqrtu-eigs}，谱中始终含有 \(\lambda_\pm(u)=\pm\sqrt u\)，而定理 \ref{thm:xi-time-part73-perron-vs-atomic-sqrtu-separation} 给出
\[
\frac{\sqrt u}{\lambda_1(u)}\longrightarrow \frac1c=\sqrt b.
\]
因此
\[
\frac{\lambda_\pm(u)}{\lambda_1(u)}\longrightarrow \pm\sqrt b,
\]
从而 \(\pm\sqrt b\in\mathcal L\)。证毕。
\end{proof}

\begin{corollary}[输运层与保护层的双层谱分裂]\label{cor:xi-time-part9w-transport-protection-two-layer-splitting}
沿用定理 \ref{thm:xi-time-part9w-zero-temp-normalized-spectrum-fourpoint-skeleton} 的记号。零温极限 \(u\to+\infty\) 中，归一化谱至少保留四点骨架
\[
\underbrace{\{1,-1\}}_{\text{输运层}}
\sqcup
\underbrace{\{\sqrt b,-\sqrt b\}}_{\text{保护层}},
\qquad 0<\sqrt b<1.
\]
其中外层对 \(\{\pm1\}\) 控制慢松弛，而内层对 \(\{\pm\sqrt b\}\) 由显式原子 Witt 因子给出。去除因子 \((1-uz^2)^{-1}\) 会抹去内层 \(\{\pm\sqrt b\}\)，但不改变首阶慢混合尺度；反之，仅保留内层而忽略镜像慢模，则无法控制 \(\tau_{\mathrm{full}}(u)\) 的主阶发散。
\end{corollary}
\begin{proof}
由 \(b=1/c^2\in(0,1)\) 可知 \(0<\sqrt b<1\)。四点同时存在由定理 \ref{thm:xi-time-part9w-zero-temp-normalized-spectrum-fourpoint-skeleton} 给出，而内层 \(\{\pm\sqrt b\}\) 正是显式特征值 \(\pm\sqrt u\) 归一化后的极限。定理 \ref{thm:xi-time-part9w-atomic-witt-surgery-slowmixing-invariance} 表明去除原子 Witt 因子后，慢混合主阶保持不变。另一方面，若只跟踪 \(\pm\sqrt b\)，则对应的归一化模始终固定为 \(\sqrt b<1\)，其关联时间尺度 \((-\log\sqrt b)^{-1}\) 保持有界；而定理 \ref{thm:xi-output-potential-endpoint-near-degeneracy-sqrtu-relaxation} 给出 \(\tau_{\mathrm{full}}(u)\asymp \sqrt u\to+\infty\)。故内层本身不能控制全谱慢松弛。证毕。
\end{proof}

\noindent
上述结果表明，本文对象中的连续 Lie 相位、\(2\)-进全息地址、Hurwitz 链 Jacobian/Prym 分裂、prime-language 共尾支撑与真实输入 \(40\) 状态核的零温谱学并不处于同一承载层级。有限二维 \(2\)-进 ambient space 足以忠实实现有限时深地址动力学，但黄金 bulk 共振既不能在任何有限维紧 Lie 相位上作为 \(L^2\) Fourier 读出出现，也不能经 dyadic solenoid 的圆因子下降为绝对连续相位密度；与此同时，零温慢混合由镜像慢模而非显式原子分支主导，而归一化谱则稳定留下 \(\{\pm1,\pm\sqrt b\}\) 的四点骨架。因而，在本文语义中，连通相位、原子保护、镜像输运与算术共振必须被视为彼此不可压平的不同层级，而任何试图把它们同时压缩到单一有限维 Lie 图景的方案都会丢失关键的证明复杂度与动力学信息。

\endinput
